% ============================================================
% Section 2: Related Work
% ============================================================

\section{Related Work}
\label{sec:related}

\subsection{Document and Image Forgery Detection}

Deep learning has substantially advanced general-purpose image forensics:
CNNs can detect subtle statistical fingerprints left by editing operations,
often outperforming hand-crafted pipelines~\citep{mantranet2019}.
\citet{cozzolino2020noiseprint} (Noiseprint) train a CNN via Siamese
contrastive learning to extract camera-model noise residuals; forgeries appear
as anomalies in this residual without requiring semantic understanding.
\citet{trufor2023} (\textsc{TruFor}) generalize this with a transformer fusion
network (NoisePrint++) that jointly models RGB semantics and noise-pattern
anomalies, achieving strong performance on general image forensics benchmarks
with ${\sim}828\,000$ pre-training images. \revB{Closely related, deepfake and
face-manipulation detection has advanced rapidly in this journal:
\citet{liang2023twostream} pair a two-stream network with hierarchical
supervision, \citet{leporoni2024guided} fuse {RGB} and depth cues, and
\citet{wang2024deepfake} add a self-attention discriminator, each improving
robustness to synthetic facial forgeries.}

Identity document forgery introduces unique challenges: forged regions are
often small (a single text field or portrait), the structured background
(guilloché patterns, UV-reactive inks) creates strong false negatives for
general-purpose noise detectors, and document type diversity demands
generalization beyond natural-image pre-training. Official evaluations on
FantasyID confirm that general-purpose baselines (TruFor, MMFusion~\citep{mmfusion2024},
FatFormer~\citep{fatformer2024}, UniFD~\citep{unifd2023}) all achieve
false-negative rates above 50\% at a 10\% false-positive
threshold~\citep{fantasyid2024}, motivating domain-specific approaches.

Among systems trained exclusively on FantasyID, \citet{edgedoc2025} propose
EdgeDoc, a lightweight CNN-Transformer encoder (EdgeNeXt~\citep{edgenext2022}) fed with the green
channel and a NoisePrint residual, with a U-Net decoder optimized jointly for
classification and localization. \citet{remtkd2025} introduce Re-MTKD, a
reinforcement-learning multi-teacher knowledge distillation framework using a
ConvNeXt~v2 encoder; the Sunlight team adapted this to identity documents using
over 60\,000 additional external images, winning both tracks of the DeepID
2025 Challenge~\citep{deepid2025challenge}.

DocVerify differs from all prior document-forensics systems in three respects:
(i) it systematically searches the task-weight $\lambda_{\text{mask}}$ via
multi-objective HPO rather than fixing it heuristically; (ii) it selects the
final configuration via a Pareto-optimal criterion; and (iii) it reports
uncertainty through nested cross-validation, providing statistically rigorous
performance estimates.

\subsection{Multi-Task Learning and HPO}

In multi-task learning, the
dominant approach combines task losses via scalar scalarization~\citep{marler2004survey}
with fixed weights set before training. \citet{kendall2018multi} propose
learning these weights via homoscedastic uncertainty; \citet{liu2019end}
introduce gradient normalization. All scalar approaches commit to a fixed
trade-off before training begins.

Multi-objective HPO extends single-objective search to jointly optimize multiple
metrics~\citep{miettinen1999nonlinear}. The Multiobjective TPE
(MOTPE)~\citep{motpe2022} adapts the TPE surrogate~\citep{tpe2011} to the
multi-objective setting, achieving competitive Pareto-front coverage with far
fewer function evaluations than evolutionary methods such as
NSGA-II~\citep{nsga2_2002}, which is critical when each evaluation requires
training a deep network. To our knowledge, this is the first work to apply Pareto-based
multi-objective HPO to joint identity-document forgery detection and
localization within a nested cross-validation framework.
